\documentclass[aspectratio=169,11pt]{beamer}

\usepackage[utf8]{inputenc}
\usepackage[T1]{fontenc}
\usepackage[ngerman]{babel}
\usepackage{lmodern}
\usepackage{tikz}
\usepackage{listings}
\usepackage{xcolor}
\usepackage{booktabs}
\usepackage{csquotes}

\usetikzlibrary{arrows.meta,calc,fit,positioning,shapes.geometric}

\definecolor{MeidanBlue}{HTML}{2563EB}
\definecolor{MeidanGreen}{HTML}{16A34A}
\definecolor{MeidanRed}{HTML}{DC2626}
\definecolor{MeidanOrange}{HTML}{F97316}
\definecolor{MeidanInk}{HTML}{111827}
\definecolor{MeidanMuted}{HTML}{6B7280}
\definecolor{MeidanPanel}{HTML}{F6F8FB}
\definecolor{Textile}{HTML}{E8DCCB}

\usetheme{Madrid}
\usecolortheme{default}
\setbeamertemplate{navigation symbols}{}
\setbeamertemplate{footline}[frame number]
\setbeamercolor{title}{fg=white,bg=MeidanBlue}
\setbeamercolor{frametitle}{fg=MeidanInk,bg=white}
\setbeamercolor{structure}{fg=MeidanBlue}
\setbeamercolor{normal text}{fg=MeidanInk,bg=white}
\setbeamercolor{block title}{fg=white,bg=MeidanBlue}
\setbeamercolor{block body}{fg=MeidanInk,bg=MeidanPanel}

\lstdefinestyle{meidan}{
  basicstyle=\ttfamily\scriptsize,
  keywordstyle=\color{MeidanBlue}\bfseries,
  stringstyle=\color{MeidanGreen},
  commentstyle=\color{MeidanMuted},
  numbers=none,
  frame=single,
  rulecolor=\color{MeidanPanel},
  breaklines=true,
  showstringspaces=false,
  columns=fullflexible
}
\lstset{style=meidan}

\lstdefinelanguage{json}{
  morestring=[b]",
  stringstyle=\color{MeidanGreen},
  literate=
   *{0}{{{\color{MeidanBlue}0}}}{1}
    {1}{{{\color{MeidanBlue}1}}}{1}
    {2}{{{\color{MeidanBlue}2}}}{1}
    {3}{{{\color{MeidanBlue}3}}}{1}
    {4}{{{\color{MeidanBlue}4}}}{1}
    {5}{{{\color{MeidanBlue}5}}}{1}
    {6}{{{\color{MeidanBlue}6}}}{1}
    {7}{{{\color{MeidanBlue}7}}}{1}
    {8}{{{\color{MeidanBlue}8}}}{1}
    {9}{{{\color{MeidanBlue}9}}}{1}
    {:}{{{\color{MeidanInk}:}}}{1}
    {,}{{{\color{MeidanInk},}}}{1}
    {\{}{{{\color{MeidanInk}\{}}}{1}
    {\}}{{{\color{MeidanInk}\}}}}{1}
    {[}{{{\color{MeidanInk}[}}}{1}
    {]}{{{\color{MeidanInk}]}}}{1},
}

\tikzset{
  component/.style={
    rounded corners=3pt,
    draw=MeidanInk!45,
    fill=white,
    line width=0.7pt,
    align=center,
    minimum width=3.1cm,
    minimum height=1.25cm,
    inner sep=6pt
  },
  smallcomponent/.style={
    rounded corners=3pt,
    draw=MeidanInk!45,
    fill=white,
    line width=0.7pt,
    align=center,
    minimum width=2.35cm,
    minimum height=0.95cm,
    inner sep=4pt
  },
  flow/.style={-Latex, line width=1pt, draw=MeidanBlue},
  altflow/.style={-Latex, line width=1pt, draw=MeidanOrange},
  textilebox/.style={
    rounded corners=4pt,
    draw=Textile!70!black,
    fill=Textile!35,
    line width=0.7pt,
    align=center,
    inner sep=8pt
  }
}

\newcommand{\apiendpoint}[1]{\texttt{#1}}

\title[Meidan]{Meidan: Interaktiver Teddybär als Baby-Wein-Alarm}
\subtitle{Prototyp für Interaktive Textilien}
\author{Projektvorstellung}
\institute{ESP32 \quad Go-Server \quad Gemini \quad Flutter-App}
\date{\today}

\begin{document}

\begin{frame}
  \titlepage
\end{frame}

\begin{frame}{Ausgangsidee}
  \small
  \begin{columns}[T,totalwidth=\textwidth]
    \begin{column}{0.55\textwidth}
      \begin{block}{Ziel}
        Ein Kuschelbär wird zu einem unauffälligen, interaktiven Textil:
        Er erkennt laute Geräusche oder mögliches Babyweinen und informiert
        die Eltern über eine App.
      \end{block}

      \vspace{0.2em}
      \begin{itemize}
        \item Elektronik verschwindet im Bär, das Objekt bleibt vertraut.
        \item Sensorwerte werden regelmäßig an einen Server gesendet.
        \item Der Server bewertet Messwerte lokal oder optional mit Gemini.
        \item Die Flutter-App zeigt Status, Pegel und Warnungen.
      \end{itemize}
    \end{column}
    \begin{column}{0.40\textwidth}
      \centering
      \begin{tikzpicture}[scale=0.85]
        \fill[Textile!65] (0,0) circle (1.25);
        \fill[Textile!75] (-0.85,1.05) circle (0.45);
        \fill[Textile!75] (0.85,1.05) circle (0.45);
        \fill[Textile!45] (-0.85,1.05) circle (0.25);
        \fill[Textile!45] (0.85,1.05) circle (0.25);
        \fill[white] (-0.45,0.25) circle (0.14);
        \fill[white] (0.45,0.25) circle (0.14);
        \fill[MeidanInk] (-0.45,0.25) circle (0.07);
        \fill[MeidanInk] (0.45,0.25) circle (0.07);
        \fill[MeidanInk] (0,-0.05) ellipse (0.16 and 0.1);
        \draw[MeidanInk, line width=0.7pt] (-0.25,-0.35) .. controls (-0.05,-0.55) and (0.05,-0.55) .. (0.25,-0.35);
        \node[smallcomponent, fill=white] at (0,-1.95) {ESP32\\Mikrofon};
        \draw[flow] (0,-1.35) -- (0,-1.45);
        \node[align=center, text=MeidanMuted] at (0,-2.85) {\small Technik im Textil\\\small statt sichtbar am Gerät};
      \end{tikzpicture}
    \end{column}
  \end{columns}
\end{frame}

\begin{frame}{Warum Interaktive Textilien?}
  \begin{columns}[T,totalwidth=\textwidth]
    \begin{column}{0.47\textwidth}
      \begin{block}{Textiles Objekt}
        \begin{itemize}
          \item weich, vertraut und kindgerecht
          \item Sensorik kann verdeckt integriert werden
          \item keine technische Hürde für die Nutzung
        \end{itemize}
      \end{block}
    \end{column}
    \begin{column}{0.47\textwidth}
      \begin{block}{Interaktion}
        \begin{itemize}
          \item Umgebung wird über Geräusche wahrgenommen
          \item Server erzeugt aus Messwerten ein Ereignis
          \item App übersetzt Ereignisse in verständliche Warnungen
        \end{itemize}
      \end{block}
    \end{column}
  \end{columns}

  \vspace{0.8em}
  \begin{center}
    \begin{tikzpicture}
      \node[textilebox, minimum width=3.2cm] (textile) {Bär / Textil};
      \node[component, right=1.2cm of textile] (sense) {Wahrnehmen\\\small Mikrofon + ESP32};
      \node[component, right=1.2cm of sense] (react) {Reagieren\\\small App-Warnung};
      \draw[flow] (textile) -- (sense);
      \draw[flow] (sense) -- (react);
    \end{tikzpicture}
  \end{center}
\end{frame}

\begin{frame}{Systemüberblick}
  \centering
  \begin{tikzpicture}[node distance=1.25cm]
    \node[textilebox, minimum width=3.2cm, minimum height=2.05cm] (bear) {
      \textbf{Teddybär}\\
      \small ESP32 + Mikrofon\\
      \small Pegel + 2s WAV
    };
    \node[component, right=1.3cm of bear] (server) {
      \textbf{Go-Server}\\
      \small REST-API\\
      \small Events + Pegel
    };
    \node[component, above right=0.4cm and 1.3cm of server] (gemini) {
      \textbf{Gemini}\\
      \small optionale\\Audioanalyse
    };
    \node[component, below right=0.4cm and 1.3cm of server] (app) {
      \textbf{Flutter-App}\\
      \small Dashboard\\
      \small Warnsignal
    };

    \draw[flow] (bear) -- node[above, font=\scriptsize]{HTTP JSON} (server);
    \draw[altflow] (server) -- node[above, sloped, font=\scriptsize]{Audio} (gemini);
    \draw[altflow] (gemini) -- node[below, sloped, font=\scriptsize]{Urteil} (server);
    \draw[flow] (app) -- node[below, sloped, font=\scriptsize]{Polling} (server);
    \draw[flow] (server) -- node[above, sloped, font=\scriptsize]{Status} (app);
  \end{tikzpicture}

  \vspace{0.8em}
  \small
  Der Server ist die zentrale Stelle: Er sammelt Messwerte, bewertet Ereignisse
  und stellt den aktuellen Zustand für die App bereit.
\end{frame}

\begin{frame}{Hardware im Bär}
  \begin{columns}[T,totalwidth=\textwidth]
    \begin{column}{0.48\textwidth}
      \begin{block}{Bauteile}
        \begin{itemize}
          \item ESP32 als Mikrocontroller mit WLAN
          \item analoges Mikrofonmodul, z. B. KY-037/KY-038
          \item Stromversorgung für den mobilen Prototyp
          \item weiche Platzierung im Inneren des Bären
        \end{itemize}
      \end{block}
    \end{column}
    \begin{column}{0.48\textwidth}
      \begin{block}{Wichtige Parameter}
        \begin{tabular}{ll}
          \toprule
          Messfenster & 120 ms\\
          Pegel & max - min\\
          Audio & 2 s, 8 kHz WAV\\
          Pegel-Upload & alle 3 s\\
          Event-Upload & bei Weinen / alle 10 s\\
          \bottomrule
        \end{tabular}
      \end{block}
    \end{column}
  \end{columns}

  \vspace{0.7em}
  \small
  Im Sketch wird der analoge Eingang \texttt{SOUND\_PIN = 34} verwendet.
  Die Empfindlichkeit wird über \texttt{CRY\_THRESHOLD} und das Potentiometer
  des Mikrofons abgestimmt.
\end{frame}

\begin{frame}[fragile]{ESP32: Messlogik}
  \begin{columns}[T,totalwidth=\textwidth]
    \begin{column}{0.50\textwidth}
      \begin{itemize}
        \item Der ESP32 liest den analogen Mikrofonwert in kurzen Fenstern.
        \item Der Pegel ist die Differenz zwischen höchstem und niedrigstem Wert.
        \item Erst nach mehreren Treffern wird ein mögliches Weinen gemeldet.
        \item Optional wird eine kurze WAV-Aufnahme als Base64 mitgesendet.
      \end{itemize}
    \end{column}
    \begin{column}{0.46\textwidth}
\begin{lstlisting}[language=C++]
SoundReading reading = readSound();
bool loudEnough =
  reading.level >= CRY_THRESHOLD;

if (loudEnough) {
  cryHitCount++;
} else {
  cryHitCount = 0;
}

bool crying =
  cryHitCount >= CRY_HITS_REQUIRED;
\end{lstlisting}
    \end{column}
  \end{columns}
\end{frame}

\begin{frame}[fragile]{HTTP-Daten}
  \begin{columns}[T,totalwidth=\textwidth]
    \begin{column}{0.47\textwidth}
      \begin{block}{Events}
        \apiendpoint{POST /api/events}
        \begin{lstlisting}[language=json]
{
  "deviceId": "esp32-1",
  "crying": true,
  "volume": 86,
  "confidence": 0.82,
  "mimeType": "audio/wav",
  "audioBase64": "..."
}
        \end{lstlisting}
      \end{block}
    \end{column}
    \begin{column}{0.47\textwidth}
      \begin{block}{Lautstärke}
        \apiendpoint{POST /api/volume}
        \begin{lstlisting}[language=json]
{
  "deviceId": "esp32-1",
  "volume": 1400
}
        \end{lstlisting}
      \end{block}
    \end{column}
  \end{columns}

  \vspace{0.4em}
  \small
  Die App liest später über \apiendpoint{GET /api/events} und
  \apiendpoint{GET /api/volume} die letzten Meldungen.
\end{frame}

\begin{frame}{Server und KI-Bewertung}
  \small
  \begin{columns}[T,totalwidth=\textwidth]
    \begin{column}{0.52\textwidth}
      \begin{block}{Aufgaben des Go-Servers}
        \begin{itemize}
          \item REST-Endpunkte für ESP32 und App
          \item Speicherung der letzten Events und Pegel im Arbeitsspeicher
          \item CORS für Flutter-Web
          \item Server-Sent Events vorbereitet über \apiendpoint{/api/events/stream}
        \end{itemize}
      \end{block}
    \end{column}
    \begin{column}{0.42\textwidth}
      \begin{block}{Gemini optional}
        \begin{itemize}
          \item Audio wird als \texttt{audio/wav} übertragen.
          \item Gemini liefert \texttt{crying}, \texttt{confidence}, \texttt{message}.
          \item Ohne API-Key oder bei Fehlern nutzt der Server lokale Regeln.
        \end{itemize}
      \end{block}
    \end{column}
  \end{columns}

  \vspace{0.2em}
  \centering
  \begin{tikzpicture}[node distance=0.75cm, scale=0.75, transform shape]
    \node[smallcomponent] (request) {ESP32\\Event};
    \node[smallcomponent, right=1.0cm of request] (server) {Go\\Server};
    \node[smallcomponent, right=1.0cm of server] (choice) {Audio?\\Key?};
    \node[smallcomponent, above right=0.35cm and 0.85cm of choice, fill=MeidanPanel] (ai) {Gemini\\Urteil};
    \node[smallcomponent, below right=0.35cm and 0.85cm of choice, fill=MeidanPanel] (local) {Lokaler\\Fallback};
    \node[smallcomponent, right=3.05cm of choice] (event) {Event\\für App};
    \draw[flow] (request) -- (server);
    \draw[flow] (server) -- (choice);
    \draw[altflow] (choice) -- (ai);
    \draw[altflow] (choice) -- (local);
    \draw[flow] (ai) -- (event);
    \draw[flow] (local) -- (event);
  \end{tikzpicture}
\end{frame}

\begin{frame}{Flutter-App für Eltern}
  \small
  \begin{columns}[T,totalwidth=\textwidth]
    \begin{column}{0.45\textwidth}
      \begin{block}{Dashboard}
        \begin{itemize}
          \item Status: Bereit, Alles ruhig, Weinen erkannt, Sehr laut
          \item aktueller Lautstärkepegel
          \item Liste der letzten Ereignisse
          \item Server-URL kann in der App gesetzt werden
        \end{itemize}
      \end{block}
    \end{column}
    \begin{column}{0.48\textwidth}
      \begin{block}{Benachrichtigung}
        \begin{itemize}
          \item App fragt Server alle 3 Sekunden ab.
          \item Neue Warnungen können einen Warnton auslösen.
          \item Testbuttons erlauben Demo ohne ESP32.
        \end{itemize}
      \end{block}
      \vspace{0.5em}
      \begin{tikzpicture}
        \node[component, minimum width=4.9cm, minimum height=2.55cm, fill=MeidanPanel] (phone) {};
        \node[anchor=north west, font=\small\bfseries, text=MeidanBlue] at ($(phone.north west)+(0.25,-0.25)$) {Meidan};
        \node[anchor=west, text=MeidanRed, font=\small\bfseries] at ($(phone.west)+(0.35,0.35)$) {Weinen erkannt};
        \node[anchor=west, text=MeidanMuted, font=\scriptsize] at ($(phone.west)+(0.35,0.0)$) {KI-Audioanalyse angefragt};
        \node[anchor=west, text=MeidanOrange, font=\small] at ($(phone.west)+(0.35,-0.55)$) {Pegel: 1400};
        \draw[MeidanBlue, line width=1pt] ($(phone.south west)+(0.35,0.45)$) -- ($(phone.south east)+(-0.35,0.45)$);
      \end{tikzpicture}
    \end{column}
  \end{columns}
\end{frame}

\begin{frame}{Datenfluss in der Demo}
  \centering
  \begin{tikzpicture}[scale=0.82, transform shape,
    flowstep/.style={smallcomponent, minimum width=2.5cm, minimum height=1.05cm},
    node distance=0.85cm
  ]
    \node[flowstep] (s1) {1\\Geräusch};
    \node[flowstep, right=of s1] (s2) {2\\ESP32 misst};
    \node[flowstep, right=of s2] (s3) {3\\HTTP an Server};
    \node[flowstep, right=of s3] (s4) {4\\Bewertung};
    \node[flowstep, right=of s4] (s5) {5\\App-Warnung};
    \draw[flow] (s1) -- (s2);
    \draw[flow] (s2) -- (s3);
    \draw[flow] (s3) -- (s4);
    \draw[flow] (s4) -- (s5);
  \end{tikzpicture}

  \vspace{1em}
  \begin{columns}[T,totalwidth=\textwidth]
    \begin{column}{0.47\textwidth}
      \begin{block}{Normaler Ablauf}
        \begin{itemize}
          \item Mikrofon erkennt niedrigen Pegel.
          \item Server speichert ruhige Meldung.
          \item App zeigt \enquote{Alles ruhig}.
        \end{itemize}
      \end{block}
    \end{column}
    \begin{column}{0.47\textwidth}
      \begin{block}{Warnfall}
        \begin{itemize}
          \item Pegel überschreitet Schwellwert oder Audio wird als Weinen bewertet.
          \item Server erzeugt Event vom Typ \texttt{cry} oder \texttt{loud}.
          \item App zeigt Warnung und kann Ton abspielen.
        \end{itemize}
      \end{block}
    \end{column}
  \end{columns}
\end{frame}

\begin{frame}[fragile]{Demo-Plan}
  \begin{columns}[T,totalwidth=\textwidth]
    \begin{column}{0.52\textwidth}
      \begin{enumerate}
        \item Go-Server starten.
        \item Flutter-Web-App starten.
        \item ESP32 im Bär einschalten.
        \item Ruhigen Zustand zeigen.
        \item Geräusch / Testevent auslösen.
        \item Warnung in der App erklären.
      \end{enumerate}
    \end{column}
    \begin{column}{0.43\textwidth}
\begin{lstlisting}[language=bash,basicstyle=\ttfamily\tiny]
cd server
go run .

cd app
flutter run -d web-server \
  --web-port=8081 \
  --dart-define=SERVER_URL=http://localhost:8080
\end{lstlisting}
    \end{column}
  \end{columns}

  \vspace{0.5em}
  \small
  Für eine Demo ohne Hardware können die Testbuttons in der App oder
  \texttt{curl}-Requests aus der README verwendet werden.
\end{frame}

\begin{frame}{Grenzen des Prototyps}
  \begin{itemize}
    \item \textbf{Audioqualität:} Ein einfaches Mikrofon im Stoff kann gedämpft oder verrauscht sein.
    \item \textbf{Kalibrierung:} Schwellwerte müssen an Raum, Mikrofon und Platzierung angepasst werden.
    \item \textbf{Datenschutz:} Audioanalyse sollte bewusst aktiviert, erklärt und abgesichert werden.
    \item \textbf{Energie:} Für einen echten Bär-Prototyp braucht es ein sicheres Akkukonzept.
    \item \textbf{Robustheit:} WLAN-Ausfälle und Serverfehler müssen für den Alltag abgefangen werden.
  \end{itemize}

  \vspace{0.9em}
  \begin{block}{Nächste Schritte}
    Gehäuse und Kabelführung im Textil verbessern, echte Push-Benachrichtigungen
    ergänzen und Messdaten mit realistischen Geräuschsituationen testen.
  \end{block}
\end{frame}

\begin{frame}{Fazit}
  \begin{columns}[T,totalwidth=\textwidth]
    \begin{column}{0.50\textwidth}
      \begin{block}{Was gezeigt wird}
        \begin{itemize}
          \item Ein Alltagsobjekt wird durch Sensorik interaktiv.
          \item ESP32, Server und App bilden eine vollständige Kette.
          \item KI kann die lokale Messlogik ergänzen, ist aber nicht zwingend nötig.
        \end{itemize}
      \end{block}
    \end{column}
    \begin{column}{0.44\textwidth}
      \begin{block}{Kernaussage}
        Interaktive Textilien müssen nicht sichtbar technisch wirken:
        Der Bär bleibt ein weiches Objekt, reagiert aber über die digitale
        Infrastruktur auf seine Umgebung.
      \end{block}
    \end{column}
  \end{columns}
\end{frame}

\appendix
\begin{frame}{Backup: API-Endpunkte}
  \begin{tabular}{lll}
    \toprule
    Methode & Pfad & Zweck\\
    \midrule
    GET & \apiendpoint{/health} & Serverstatus\\
    GET & \apiendpoint{/api/events} & letzte Events\\
    POST & \apiendpoint{/api/events} & Weinen-/Ruhe-Event anlegen\\
    GET & \apiendpoint{/api/events/stream} & Event-Stream\\
    GET & \apiendpoint{/api/volume} & letzte Pegelwerte\\
    POST & \apiendpoint{/api/volume} & Pegelwert anlegen\\
    \bottomrule
  \end{tabular}
\end{frame}

\end{document}
